

\section{Techniques d'ancrage moteur : TPR et Living Sequential Expression}
space{0.3cm}
oindent



\subsection{Point iii.2.2 : total physical response (tpr) et living sequential expression (lse) pour l'inscription en mémoire procédurale}



\subsection{Sommaire exécutif}

Ce rapport constitue une analyse approfondie et critique du point III.2.\cite{III-2-2-ref2} de votre
plan de recherche, focalisé sur l'utilisation des techniques d'ancrage moteur — spécifiquement la
\textit{Total Physical Response} (TPR) et le \textit{Living Sequential Expression} (LSE) — dans le
cadre de la didactique du grec ancien. L'objectif central de cette investigation est d'évaluer
comment ces méthodologies, initialement conçues pour les langues vivantes modernes, peuvent être
adaptées pour revitaliser l'enseignement d'une langue de corpus, en ciblant spécifiquement la
mémoire procédurale (implicite) plutôt que la mémoire déclarative (explicite) traditionnellement
sollicitée par la méthode grammaire-traduction.

L'analyse démontre que l'approche traditionnelle, focalisée sur l'analyse morphosyntaxique
abstraite, échoue souvent à produire une compétence de lecture fluide car elle sature la mémoire de
travail sans automatiser les processus linguistiques de bas niveau. À l'inverse, la TPR et la LSE
exploitent les mécanismes de la cognition incarnée (\textit{embodied cognition}) et de l'effet
d'enactment (\textit{enactment effect}) pour créer des traces mnésiques robustes et
multisensorielles. En liant le signe linguistique à l'exécution motrice et à la séquence narrative,
ces méthodes permettent de \og hacker\fg{} les systèmes neurobiologiques d'acquisition du langage,
transformant le grec ancien d'un objet d'étude statique en une compétence dynamique. Le rapport
détaille les fondements neuroscientifiques de ces approches, leur historique pédagogique depuis
François Gouin jusqu'à l'Institut Polis, et propose une synthèse opérationnelle pour leur
intégration curriculaire.

L'enseignement des langues classiques, et du grec ancien en particulier, traverse une crise de
légitimité et d'efficacité pédagogique qui perdure depuis la fin du XXe siècle. Traditionnellement,
l'apprentissage du grec repose sur le \og modèle prussien\fg{} ou méthode grammaire-traduction, qui
privilégie la mémorisation par cœur de paradigmes morphologiques décontextualisés et l'analyse
syntaxique de textes littéraires complexes dès les premiers stades
d'apprentissage.\cite{III-2-2-ref1} Cette approche, bien qu'intellectuellement rigoureuse, produit
souvent ce que les chercheurs qualifient de \og connaissance inerte\fg{} : l'étudiant est capable de
réciter la déclinaison de \textit{logos} ou la conjugaison de l'aoriste passif, mais se trouve dans
l'incapacité de comprendre une phrase simple sans recourir à une analyse consciente et laborieuse,
médiée par la langue maternelle.

Le point III.2.\cite{III-2-2-ref2} de votre plan de recherche intervient à un moment charnière où la
didactique des langues anciennes cherche à intégrer les acquis des sciences cognitives et de la
linguistique appliquée. L'hypothèse centrale que nous explorerons est que le déficit de fluidité
observé chez les étudiants en grec ancien ne provient pas d'un manque d'effort intellectuel, mais
d'une erreur de ciblage cognitif : l'enseignement traditionnel s'adresse presque exclusivement à la
\textbf{mémoire déclarative}, alors que la maîtrise d'une langue — même ancienne — requiert
l'engagement massif de la \textbf{mémoire procédurale}.\cite{III-2-2-ref3}

Pour remédier à cette asymétrie, les techniques d'ancrage moteur émergent comme des solutions
privilégiées. Elles reposent sur le postulat que le langage n'est pas une abstraction logique
flottant dans l'esprit, mais une activité biologique enracinée dans l'expérience sensori-motrice du
corps. Ce rapport se propose d'examiner deux incarnations majeures de ce principe : la \textit{Total
Physical Response} (TPR), qui lie l'impératif au mouvement immédiat, et le \textit{Living Sequential
Expression} (LSE), qui lie la narration à la séquence d'actions logiques. Nous verrons comment ces
méthodes, loin d'être de simples artifices ludiques, constituent des leviers neurobiologiques
puissants pour graver la structure du grec ancien dans les réseaux neuronaux profonds de
l'apprenant.



\subsection{Fondements neurocognitifs : du mouvement à la mémoire}

Pour comprendre la pertinence de la TPR et de la LSE, il est impératif de disséquer les mécanismes
cérébraux qu'elles sollicitent. L'idée que \og faire\fg{} aide à \og retenir\fg{} n'est pas
nouvelle, mais les neurosciences contemporaines offrent désormais un cadre explicatif robuste,
notamment à travers le modèle Déclaratif/Procédural et la théorie de la cognition incarnée.



\subsection{Le modèle déclaratif/procédural (dp) de michael ullman}

Le cadre théorique le plus pertinent pour notre analyse est le \textbf{modèle déclaratif/procédural
(DP)} développé par le neuroscientifique Michael Ullman. Ce modèle postule que le langage dépend de
deux systèmes de mémoire et d'apprentissage distincts, qui sont également responsables d'autres
fonctions cognitives non linguistiques.\cite{III-2-2-ref3}

La Mémoire Déclarative : Le Savoir Explicite

Ce système, ancré neuro-anatomiquement dans le lobe temporal médian (incluant l'hippocampe) et le
néocortex, est spécialisé dans l'apprentissage, le stockage et la récupération de faits et
d'événements. C'est une mémoire explicite, accessible à la conscience.

\begin{itemize}\item \textit{Dans l'apprentissage du grec :} C'est le système utilisé pour apprendre le vocabulaire (le sens arbitraire des mots) et les règles de grammaire explicites (\og l'aoriste indique une action ponctuelle\fg{}).\item \textit{Limitation :} Bien que rapide pour l'encodage initial, la récupération d'informations dans la mémoire déclarative est cognitivement coûteuse et lente. Elle nécessite un effort attentionnel constant.\cite{III-2-2-ref4}\end{itemize}La Mémoire Procédurale : Le Savoir-Faire Implicite

Ce système est enraciné dans les circuits fronto-striataux, comprenant les ganglions de la base
(notamment le noyau caudé), le cortex prémoteur et l'aire de Broca.\cite{III-2-2-ref7} Il est
spécialisé dans l'apprentissage et le contrôle des habiletés motrices et cognitives séquentielles,
comme faire du vélo ou jouer d'un instrument.

\begin{itemize}\item \textit{Dans l'apprentissage du grec :} Pour un locuteur compétent, c'est ce système qui gère la grammaire (la combinatoire des mots) de manière rapide, automatique et non consciente. C'est la \og mémoire du muscle\fg{} appliquée à la syntaxe.\cite{III-2-2-ref3}\end{itemize}L'Enjeu Pédagogique : La \og Proceduralization\fg{}

Le problème fondamental de la méthode grammaire-traduction est qu'elle confine l'apprenant à la
mémoire déclarative. L'étudiant sait comment construire un optatif, mais il ne possède pas la
procédure pour le générer ou le comprendre en temps réel. Les techniques d'ancrage moteur comme la
TPR et la LSE visent à faciliter le basculement (shift) ou l'interaction entre ces deux systèmes. En
liant la langue à une séquence motrice (gérée par le système procédural), on force le cerveau à
traiter l'information linguistique via les circuits striataux, favorisant ainsi
l'automatisation.\cite{III-2-2-ref6}



\subsection{Cognition incarnée et effet d'enactment}

Au-delà de la distinction entre les systèmes de mémoire, la théorie de la \textbf{Cognition
Incarnée} (\textit{Embodied Cognition}) remet en cause la vision computationnelle classique du
cerveau. Selon cette approche, la compréhension du langage implique une simulation mentale des
actions décrites.

La Simulation Neurale

Lorsque nous entendons le verbe \og courir\fg{} (trechein en grec), notre cerveau n'accède pas
seulement à une définition abstraite dans un dictionnaire mental. Il réactive, à un niveau
subliminaire, les zones du cortex moteur responsables du mouvement des jambes.\cite{III-2-2-ref10}
Les recherches montrent que le traitement des verbes d'action active de manière somatotopique le
cortex moteur et prémoteur : lire \og saisir\fg{} active la zone de la main, lire \og botter\fg{}
active la zone du pied.\cite{III-2-2-ref13}

L'Effet d'Enactment (Enactment Effect)

Cette connexion neurale explique l'Effet d'Enactment (ou effet de la tâche réalisée par le sujet),
abondamment documenté par des chercheurs comme Engelkamp, Zimmer et Macedonia.\cite{III-2-2-ref9}
Les études démontrent de manière cohérente que les mots et expressions appris en étant physiquement
exécutés (encodage moteur) sont mémorisés de manière significativement supérieure à ceux appris par
simple écoute ou répétition verbale.

\begin{itemize}\item \textit{Mécanisme :} L'exécution du geste crée une trace mnésique multimodale complexe. L'information n'est pas seulement stockée sous forme phonologique (le son) ou sémantique (le sens), mais aussi sous forme sensorimotrice et proprioceptive.\cite{III-2-2-ref9} Cette redondance de l'encodage rend la trace mnésique plus stable et plus facile à récupérer.\item \textit{Trace Theory :} Comme le soulignait James Asher en se référant à la théorie de la trace de Katona (1940), plus une connexion est \og tracée\fg{} intensément (par la répétition verbale ET motrice), plus l'association est forte.\cite{III-2-2-ref16}\end{itemize}

\subsection{Implications pour les langues "mortes"}

L'application de ces principes au grec ancien est révolutionnaire car elle propose de traiter une
langue \og morte\fg{} (sans locuteurs natifs disponibles pour l'immersion sociale) comme une langue
\og vivante\fg{} pour le cerveau. Le cerveau ne distingue pas si la langue est parlée dans la rue ou
non ; il réagit aux stimuli qu'on lui présente. Si l'enseignement du grec ancien utilise l'ancrage
moteur, il active les mêmes mécanismes de plasticité cérébrale que pour l'apprentissage d'une langue
vivante moderne.\cite{III-2-2-ref18} La TPR et la LSE ne sont donc pas des simulations
artificielles, mais des stimulateurs réels de la neurobiologie du langage.



\subsection{Total physical response (tpr) : l'impératif comme levier d'acquisition}

La méthode \textbf{Total Physical Response (TPR)}, développée dans les années 1960 et 1970 par le Dr
James Asher, professeur de psychologie à l'Université d'État de San José, représente la première
tentative systématique d'aligner l'enseignement des langues sur les principes de l'apprentissage
moteur et du développement de l'enfant.\cite{III-2-2-ref20} Pour l'enseignant de grec ancien, elle
offre une porte d'entrée pragmatique pour construire les fondations de la langue sans passer par
l'abstraction grammaticale.



\subsection{Les hypothèses fondatrices de james asher}

La TPR n'est pas une simple collection de jeux de classe ; elle repose sur une architecture
théorique précise, articulée autour de trois hypothèses majeures que l'enseignant de grec doit
comprendre pour appliquer la méthode efficacement.\cite{III-2-2-ref22}



\subsection{A. l'hypothèse du bio-programme inné (\textit{the bio-program})}

Asher part du constat que l'apprentissage d'une langue seconde (L2) est souvent un échec laborieux,
alors que l'acquisition de la langue maternelle (L1) est un succès universel et apparemment sans
effort. Il en déduit l'existence d'un \og bio-programme\fg{} inné qui définit la séquence optimale
d'apprentissage.\cite{III-2-2-ref16}

1. \textbf{Priorité à la Compréhension :} Dans le développement naturel, la compétence d'écoute
précède toujours la compétence de parole. Les enfants passent une longue \og période
silencieuse\fg{} à décoder le flux sonore avant de produire leurs premiers mots.

2. \textbf{L'Impératif comme Vecteur :} Une grande partie de l'input linguistique reçu par l'enfant
prend la forme de commandes parentales (\og Tiens ça\fg{}, \og Viens ici\fg{}, \og Regarde\fg{}). La
compréhension est validée par l'action physique de l'enfant, pas par une réponse verbale.

3. \textbf{Émergence Naturelle de la Parole :} La production orale ne doit pas être forcée. Elle
émerge spontanément une fois que la carte cognitive de la langue est suffisamment établie par
l'écoute et l'action.

\textit{Application au Grec :} Cela signifie que les premiers mois d'apprentissage du grec ne
devraient pas exiger de l'étudiant qu'il traduise ou qu'il construise des phrases, mais qu'il
réagisse physiquement à des instructions en grec.



\subsection{B. l'hypothèse de la latéralisation cérébrale}

Bien que la neuroscience ait évolué depuis les années 1970, l'intuition d'Asher reste pertinente sur
le plan fonctionnel. Il postulait que l'apprentissage traditionnel (analyse, règles) sollicite
l'hémisphère gauche, tandis que l'apprentissage moteur engage l'hémisphère
droit.\cite{III-2-2-ref22} Asher suggérait que l'hémisphère droit peut traiter le langage de manière
globale et intuitive par le mouvement, contournant ainsi la \og résistance\fg{} analytique de
l'hémisphère gauche. Pour une langue à la morphologie aussi complexe que le grec, passer par
l'intuition motrice permet d'éviter la paralysie de l'analyse (\og Est-ce un accusatif ou un
génitif?\fg{}) au profit de la réaction (\og Je dois aller vers la porte\fg{}).



\subsection{La réduction du stress (le filtre affectif)}

S'inspirant de la psychologie humaniste, Asher insiste sur le fait que le stress inhibe
l'apprentissage (le \og filtre affectif\fg{} de Krashen).\cite{III-2-2-ref16} En TPR, l'apprenant
n'est jamais mis en échec public. S'il ne comprend pas un ordre, il lui suffit d'observer et
d'imiter ses camarades ou l'enseignant. Cette sécurité psychologique est cruciale pour le grec
ancien, souvent perçu comme une langue élitiste et intimidante où l'erreur est sévèrement
sanctionnée.



\subsection{Méthodologie opérationnelle en classe de grec}

La mise en œuvre de la TPR en cours de grec koinè ou classique suit une progression rigoureuse,
utilisant l'\textbf{impératif} comme noyau grammatical central. Contrairement aux idées reçues,
l'impératif permet d'enseigner bien plus que de simples ordres ; il permet d'ancrer la structure
casuelle et prépositionnelle.

\begin{table}[h!]\small\centering\begin{tabular}{| l | l | l | l | l |}\hline\textbf{Phase} & \textbf{Description de l'Activité} & \textbf{Exemple en Grec (Enseignant)} & \textbf{Action de l'Apprenant (Sans parole)} & \textbf{Insight Grammatical (Implicite)} \\ \hline\textbf{1. Modélisation} & L'enseignant énonce l'ordre et l'exécute lui-même. & \textit{Ἀνάστηθι\!} (Lève-toi) / \textit{Κάθιζε\!} (Assieds-toi) & Observe et associe le son au geste. & Verbes de mouvement de base. Distinction singulier/pluriel (\textit{\-thi} vs \textit{\-te}). \\ \hline\textbf{2. Pratique Guidée} & L'enseignant donne l'ordre à la classe ou à des volontaires. & \textit{Bαῖνε\!} (Marche) / \textit{Στῆθι\!} (Arrête-toi) & Exécute l'action. & Validation immédiate de la compréhension. \\ \hline\textbf{3. Expansion (Objets)} & Introduction d'objets directs (Accusatif). & \textit{Ἄρον τὸ βιβλίον.} (Prends le livre) & Prend le livre. & L'Accusatif (\textit{\-on}) est perçu comme l'objet manipulé. \\ \hline\textbf{4. Expansion (Espace)} & Introduction de prépositions spatiales. & \textit{Bαῖνε πρὸς τὴν θύραν.} (Marche vers la porte) & Marche vers la porte. & L'association \textit{pros} \+ Accusatif devient une sensation de direction/mouvement. \\ \hline\textbf{5. Nouveauté/Recombinaison} & Combinaison inédite de mots connus. & \textit{Ἄρον τὸ βιβλίον καὶ βάλε ἐπὶ τὴν κεφαλήν.} (Prends le livre et mets-le sur la tête) & Exécute la séquence complexe. & Compréhension de la syntaxe et des connecteurs (\textit{kai}). \\ \hline\end{tabular}\end{table}Le Traitement de la Morphologie Verbale :

Une critique courante est que la TPR ne permet pas d'enseigner les subtilités des temps grecs. Or,
Asher et ses successeurs (comme Christophe Rico) ont montré que l'aspect verbal peut être
mimé.\cite{III-2-2-ref26}

\begin{itemize}\item \textbf{Aspect Présent (Duratif) :} L'enseignant utilise l'impératif présent (\textit{Bαῖνε} \- Continue de marcher) pour une action qui dure.\item Aspect Aoriste (Ponctuel) : L'enseignant utilise l'impératif aoriste (Bῆσον \- Fais un pas / Mets-toi en marche) pour une action ponctuelle ou ingressive.\end{itemize}L'étudiant \og sent\fg{} physiquement la différence entre le processus (Présent) et l'événement
(Aoriste) bien avant d'apprendre les règles de formation du thème verbal.



\subsection{Efficacité et limites : ce que dit la recherche}

Les études comparatives sur l'efficacité de la TPR, bien que souvent menées sur des langues vivantes
(anglais, espagnol), offrent des résultats transposables. Les recherches citées dans les snippets 20
indiquent que :

\begin{itemize}\item La TPR produit une rétention lexicale à long terme significativement supérieure à la méthode grammaire-traduction.\item Elle est particulièrement efficace pour les débutants (niveaux A1-A2) et pour le vocabulaire concret.\item Elle réduit l'attrition (abandon) des étudiants, un problème majeur dans les cursus de langues anciennes.\cite{III-2-2-ref20}\end{itemize}Cependant, la TPR présente des limites intrinsèques pour le grec ancien 29 :

\begin{itemize}\item \textbf{Le Plafond de l'Abstraction :} Il est difficile de mimer des concepts comme \og la vertu\fg{} (\textit{aretē}) ou \og la justice\fg{} (\textit{dikaiosunē}) sans tomber dans des charades complexes.\item La Passivité Linguistique : Si elle est utilisée exclusivement, la TPR peut enfermer l'étudiant dans un rôle d'exécutant muet. Elle ne prépare pas directement à la production de phrases complexes ou à la narration.\end{itemize}C'est ici que la seconde technique, le Living Sequential Expression, devient indispensable pour
prendre le relais.



\subsection{Living sequential expression (lse) : de la séquence à la narration}

Si la TPR pose les fondations lexicales et instinctives, le \textbf{Living Sequential Expression
(LSE)} construit l'architecture narrative et syntaxique. Cette technique, remise au goût du jour par
l'Institut Polis à Jérusalem, trouve ses racines dans une intuition pédagogique du XIXe siècle aussi
géniale qu'oubliée : la \og Méthode des Séries\fg{} de François Gouin.



\subsection{Archéologie de la méthode : l'épiphanie de françois gouin}

L'histoire de la LSE commence par un échec retentissant, relaté dans l'ouvrage de François Gouin,
\textit{L'Art d'enseigner et d'étudier les langues} (1880). Professeur de latin, Gouin tente
d'apprendre l'allemand en appliquant rigoureusement les méthodes académiques de son temps : il
mémorise une grammaire entière et les 248 verbes irréguliers en dix jours. Arrivé à l'université de
Berlin, il découvre avec horreur qu'il est incapable de comprendre un seul mot des cours. Il tente
alors de mémoriser le dictionnaire entier (30 000 mots en un mois), mais le résultat est le même :
\og Mes yeux voyaient les mots allemands, mais mes oreilles restaient
fermées\fg{}.\cite{III-2-2-ref30}

De retour en France, épuisé et convaincu de son inaptitude, il observe son neveu de trois ans qui,
après une visite au moulin, rejoue la scène en se parlant à lui-même. L'enfant ne récite pas des
mots isolés (\og farine\fg{}, \og meunier\fg{}), mais décompose l'événement complexe en une chaîne
causale d'actions simples : \og Il prend le sac, il va au moulin, il verse le grain, le grain tombe,
la roue tourne...\fg{}.\cite{III-2-2-ref33}

Gouin tire de cette observation trois principes révolutionnaires qui fondent la \textbf{Méthode des
Séries} :

1. \textbf{L'Organisation Séquentielle :} Le cerveau humain mémorise les événements (et donc le
langage qui les décrit) selon un ordre chronologique et logique strict. Apprendre des mots hors
contexte est contre-nature.

2. \textbf{L'Incubation Mentale :} L'enfant visualise l'action mentalement avant de la verbaliser.
Le langage sert à \og transformer des perceptions en conceptions\fg{}.

3. \textbf{Le Verbe comme Noyau :} Le verbe est l'âme de la phrase ; c'est lui qui porte la
séquence. Gouin développe des séries de 15 à 20 phrases décrivant chaque aspect de la vie (la
maison, le feu, l'eau), où chaque phrase contient un verbe d'action spécifique.\cite{III-2-2-ref35}

L'Exemple Canonique : La Série de la Porte

Pour illustrer sa méthode, Gouin a créé une série célèbre décomposant l'action d'ouvrir une porte.
Cette série est citée dans presque toutes les sources historiques et constitue le modèle de base
pour l'adaptation au grec 33 :

\begin{itemize}\item Je marche vers la porte. (\textit{I walk towards the door})\item Je m'approche de la porte. (\textit{I draw near to the door})\item J'arrive à la porte. (\textit{I get to the door})\item Je m'arrête à la porte. (\textit{I stop at the door})\item J'étends le bras. (\textit{I stretch out my arm})\item Je saisis la poignée. (\textit{I take hold of the handle})\item Je tourne la poignée. (\textit{I turn the handle})\item J'ouvre la porte. (\textit{I open the door})\item Je tire la porte. (\textit{I pull the door})\item La porte tourne sur ses gonds. (\textit{The door turns on its hinges})\item Je lâche la poignée. (\textit{I let go of the handle})\end{itemize}

\subsection{Le renouveau à l'institut polis : la technique lse}

Christophe Rico, linguiste et doyen de l'Institut Polis, a formalisé l'héritage de Gouin sous le nom
de \textbf{Living Sequential Expression (LSE)}, en l'adaptant spécifiquement aux défis des langues
anciennes (Grec Koinè, Hébreu, Latin).\cite{III-2-2-ref19}

Contrairement à la simple mémorisation de textes (comme dans certaines approches \og naturelles\fg{}
du début du XXe siècle), la LSE chez Polis est une technique active qui combine la TPR et la
production orale. Elle vise à faire passer l'étudiant du statut de récepteur d'ordres (TPR) à celui
de narrateur (LSE).

La Mécanique du LSE : La Triangulation des Personnes

L'innovation majeure de Rico est l'introduction systématique de la conjugaison par le biais d'un
dialogue triangulaire pendant l'action. Une séquence LSE typique en cours de grec se déroule ainsi
26 :

Étape 1 : L'Ordre (Impératif \- TPR)

L'enseignant donne un ordre à l'étudiant A.

\begin{itemize}\item Enseignant : \textit{Ἀνάτρεψον τὴν τράπεζαν\!} (Renverse la table\! \- \textit{verbe à l'impératif aoriste})\item Étudiant A : Exécute l'action physiquement.\end{itemize}Étape 2 : L'Auto-Narration (1ère Personne)

Pendant ou juste après l'action, l'enseignant interroge l'étudiant A.

\begin{itemize}\item Enseignant : \textit{Τί ἐποίησας;} (Qu'as-tu fait?)\item Étudiant A : \textit{Ἀνέτρεψα τὴν τράπεζαν.} (J'ai renversé la table. \- \textit{Indicatif aoriste, 1ère pers.})\item \textit{Ancrage Moteur :} L'étudiant associe la terminaison \textit{\-sa} (aoriste sigmatique 1ère pers.) à son propre geste et à son intention (\og Je\fg{}).\end{itemize}Étape 3 : L'Hétéro-Narration (3ème Personne)

L'enseignant se tourne vers l'étudiant B ou le reste de la classe.

\begin{itemize}\item Enseignant : \textit{Τί ἐποίησε;} (Qu'a-t-il fait?)\item Étudiant B : \textit{Ἀνέτρεψε τὴν τράπεζαν.} (Il a renversé la table. \- \textit{Indicatif aoriste, 3ème pers.})\item \textit{Ancrage Moteur :} La terminaison \textit{\-se} est associée à l'observation de l'action d'autrui.\end{itemize}Étape 4 : La Séquence Complète (Gouin Series)

Une fois les verbes individuels maîtrisés par ce processus, l'enseignant enchaîne les actions pour
former une histoire logique (la \og Série\fg{}). Les étudiants doivent ensuite raconter toute la
séquence, ce qui les oblige à utiliser des connecteurs logiques (kai, eita, tote).



\subsection{Exemples d'application au corpus grec}

L'Institut Polis a développé des manuels et du matériel pédagogique basés sur ces séquences,
illustrant des actions de la vie quotidienne antique ou moderne.\cite{III-2-2-ref19}

Série \og L'Écriture\fg{} (Adaptation pour le cours de grec)

Cette série permet d'introduire le vocabulaire de l'école et de l'écriture, omniprésent dans les
textes (papyrus, calame, écrire).

1. \textit{Zητῶ τὸν κάλαμον.} (Je cherche le calame/roseau.)

2. \textit{Eὑρίσκω τὸν κάλαμον.} (Je trouve le calame.)

3. \textit{Λαμβάνω τὸν κάλαμον τῇ χειρί.} (Je prends le calame avec la main. \- \textit{Introduction
du Datif instrumental})

4. \textit{Bάπτω τὸν κάλαμον εἰς τὸ μέλαν.} (Je trempe le calame dans l'encre.)

5. \textit{Γράφω τὰ γράμματα ἐπὶ τοῦ χάρτου.} (J'écris les lettres sur le papier/papyrus.)

6. \textit{Ἀποτίθημι τὸν κάλαμον.} (Je dépose le calame.)

Bénéfice Syntaxique :

Dans cette séquence, l'étudiant pratique la différence entre les prépositions de mouvement (eis \+
Accusatif : dans l'encre) et de lieu (epi \+ Génitif : sur le papier). La \og règle\fg{} n'est
jamais expliquée abstraitement ; elle est vécue. La main qui trempe \og dans\fg{} (eis) fait un
mouvement différent de la main qui écrit \og sur\fg{} (epi). La mémoire procédurale encode cette
distinction spatiale motrice.



\subsection{Synthèse et synergies : vers une pédagogie intégrée}

L'analyse séparée de la TPR et de la LSE révèle leurs forces respectives, mais c'est leur
articulation au sein d'un curriculum cohérent qui permet de véritablement \og graver\fg{} le grec
ancien dans la mémoire procédurale.



\subsection{Comparaison stratégique}

Le tableau ci-dessous synthétise les différences fonctionnelles et la complémentarité des deux
approches.

\begin{table}[h!]\small\centering\begin{tabular}{| p{2.3cm} | p{3.5cm} | p{3.5cm} |}\hline\textbf{Dimension} & \textbf{Total Physical Response (TPR)} & \textbf{Living Sequential Expression (LSE)} \\ \hline\textbf{Origine Historique} & Psychologie comportementale \& humaniste (Asher, 1960s) & Linguistique naturaliste (Gouin, 1880s) / Rico (2010s) \\ \hline\textbf{Focalisation Cognitive} & Réception (Input), Compréhension, Réaction réflexe & Production (Output), Narration, Logique causale \\ \hline\textbf{Mode Verbal Dominant} & Impératif (\textit{Λαβέ}) & Indicatif (\textit{Λαμβάνω}, \textit{Λαμβάνει}) \\ \hline\textbf{Rôle de l'Apprenant} & Exécutant physique (silencieux au début) & Acteur-Narrateur (parle en agissant) \\ \hline\textbf{Apport Grammatical} & Vocabulaire concret, Prépositions spatiales, Morphologie de l'impératif & Conjugaison (Personnes, Temps), Connecteurs logiques, Structure de phrase complexe \\ \hline\textbf{Type de Mémoire} & Association Son-Geste immédiate (Réflexe) & Association Séquence-Sens (Reconstruction mentale) \\ \hline\end{tabular}\end{table}

\subsection{Progression curriculaire suggérée}

Pour un cours de grec ancien visant la fluidité (comme le programme de Master en Philologie Ancienne
de l'Institut Polis), l'intégration se fait par étapes successives 18 :

1. Phase d'Immersion Initiale (Semaines 1-4) : Dominance TPR

L'objectif est de briser le filtre affectif et d'installer les 100-200 premiers mots (verbes de
mouvement, objets de la classe, corps humain). L'enseignant parle, les élèves agissent. Aucune
traduction n'est permise. La mémoire phonologique se calibre.

2. Phase de Transition (Semaines 5-10) : Introduction du LSE Simple

Dès que le vocabulaire de base est acquis, l'enseignant introduit la question Ti poieis? (Que fais-
tu?). On passe de l'ordre à la description. C'est le moment critique où la morphologie verbale (les
terminaisons \-ō, \-eis, \-ei) est connectée à l'identité de l'agent (Je, Tu, Il).

3. Phase de Consolidation (Semaines 11+) : LSE Complexe et Narratif

Les élèves apprennent des séries entières (Gouin Series) par cœur, en les jouant. Cela leur donne
des \og blocs\fg{} de langue préfabriqués (des chunks procéduraux) qu'ils peuvent ensuite recombiner
pour raconter des histoires plus libres ou résumer des mythes simples (ex: Les Travaux d'Hercule
traités comme une série d'actions).

4. Phase Textuelle : Le Retour au Corpus

Une fois ces automatismes procéduraux en place, la lecture des textes originaux change de nature.
Lorsqu'un étudiant lit Exelthen oiesous (\og Jésus sortit\fg{}), son cerveau ne cherche pas la règle
de l'aoriste second ; il active la trace motrice du verbe exerchomai (sortir) qu'il a physiquement
mimé des centaines de fois en classe via la TPR et la LSE. La lecture devient une simulation interne
fluide.



\subsection{Le défi de l'abstraction et la "métaphore motrice"}

Une critique récurrente concerne la capacité de ces méthodes à enseigner le vocabulaire abstrait de
la philosophie ou de la théologie.\cite{III-2-2-ref29} Cependant, la linguistique cognitive (Lakoff
\& Johnson) nous enseigne que la plupart des concepts abstraits sont des métaphores spatiales ou
motrices.

\begin{itemize}\item \textit{Comprendre} (\textit{suniēmi}) en grec signifie littéralement \og envoyer ensemble/rassembler\fg{}.\item Péché (hamartia) signifie \og manquer la cible\fg{} (terme de tir à l'arc).\end{itemize}En LSE, on peut enseigner le sens littéral/physique de ces mots par une série (tirer à l'arc,
manquer la cible), créant ainsi une base sensorimotrice solide sur laquelle le sens théologique
abstrait peut se greffer. Le cerveau traite alors l'abstraction comme une extension du concret, et
non comme un pur symbole arbitraire.



\subsection{Conclusion}

L'intégration des techniques d'ancrage moteur, TPR et LSE, dans l'enseignement du grec ancien ne
relève pas d'une mode pédagogique passagère, mais d'un réalignement nécessaire de la didactique sur
la réalité neurobiologique de l'apprentissage humain.

L'analyse des sources et des mécanismes cognitifs permet de tirer les conclusions suivantes :

1. \textbf{La Primauté du Procédural :} La fluidité en lecture, objectif ultime des études
classiques, dépend de l'automatisation des processus de bas niveau, gérée par la mémoire
procédurale. La méthode grammaire-traduction échoue souvent à atteindre ce stade car elle ne
sollicite pas les circuits fronto-striataux liés à l'action.

2. \textbf{La Complémentarité TPR/LSE :} La TPR est l'outil idéal pour l'initialisation (création de
la carte phonologique et lexicale sans stress), tandis que la LSE est l'outil de la structuration
syntaxique et narrative. Ensemble, elles couvrent le spectre de l'acquisition, du mot isolé au
discours complexe.

3. \textbf{La Revitalisation du \og Mort\fg{} :} En simulant une immersion par l'action et la
narration séquentielle, ces méthodes créent une \og mémoire procédurale artificielle\fg{} pour le
grec ancien. Elles permettent à l'étudiant de vivre la langue de l'intérieur, transformant le
déchiffrement laborieux en une expérience de compréhension directe et incarnée.

Pour l'enseignant de grec ancien au XXIe siècle, adopter ces techniques demande un changement de
paradigme : il ne s'agit plus de transmettre un savoir sur la langue (métalinguistique), mais de
guider une expérience de la langue. C'est à ce prix que les textes d'Homère, de Platon et du Nouveau
Testament pourront, à nouveau, parler directement à l'esprit des étudiants.

